\documentclass{article}
\usepackage[UTF8]{ctex}
\usepackage{amsmath,amssymb,amsfonts}
\usepackage{enumitem}
\usepackage{geometry}
\geometry{a4paper, margin=1in}
% === 错题本排版优化（不参与题目标识） ===
\setlength{\parindent}{0pt}
\setlength{\parskip}{0.35em plus 0.10em minus 0.08em}
\setlength{\emergencystretch}{3em}
\setlist{itemsep=0.25em, topsep=0.40em, parsep=0pt}
\allowdisplaybreaks[2]
% === 错题本排版优化结束 ===
\begin{document}
\section*{中国科学院大学2025年高等代数真题}
\section*{14．中国科学院大学2025年高等代数试卷}

\subsection{题目1}
1．（18 分）设 \(F(x)=\prod_{i=1}^{n}\left(x-a_{i}\right)\) ，其中 \(a_{1}, a_{2}, \cdots, a_{n}\) 是两两不同的实数。

\begin{itemize}
  \item[（1）] 记 \(l_{i}(x)=\frac{F(x)}{\left(x-a_{i}\right) F^{\prime}\left(a_{i}\right)}\) ，证明：\(l_{1}(x), l_{2}(x), \cdots, l_{n}(x)\) 线性无关。
  \item[（2）] 证明：\(\sum_{i=1}^{n} l_{i}(x)=1\) 。
  \item[（3）] 任取次数小于 \(n\) 的多项式 \(f(x)\) ，将 \(f(x)\) 用 \(l_{1}(x), l_{2}(x), \cdots, l_{n}(x)\) 进行线性表示．
\end{itemize}

\subsection{题目2}
2．（20 分）设 \(n\) 阶矩阵 \(A\) 为实对称正定矩阵 \((n>2)\) 。证明：\(A=I\) 当且仅当 \(A^{*}=A\) 。

\subsection{题目3}
3．（20 分）设 \(n+1\) 阶矩阵 \(A\) 有特征值 0 和 \(n\) 个不同的单位根 \(\omega_{1}, \omega_{2}, \cdots, \omega_{n}\left(\left(\omega_{i}\right)^{n}=1\right)\) 。证明：\(2 I-A\) 可逆，并求其逆矩阵。

\subsection{题目4}
4．（20 分）设 \(B \in M^{n \times n}(\mathbb{R})\) 为正定矩阵，\(C \in M^{n \times m}(\mathbb{R})\) 为列满秩矩阵，记

\[
A=\left(\begin{array}{cc}
B & C \\
C^{\mathrm{T}} & O
\end{array}\right) .
\]

求 \(A\) 的正负惯性指数。

\subsection{题目5}
5．（20 分）数域 \(\mathbb{F}\) 上的 \(n\) 阶矩阵 \(A\) 有 \(n\) 个不同的特征值 \(\lambda_{1}, \lambda_{2}, \cdots, \lambda_{n}, \alpha_{1}, \alpha_{2}, \cdots, \alpha_{n}\) 为对应的特征向量，记 \(\beta=\alpha_{1}+\alpha_{2}+\cdots+\alpha_{n}\) 。

\begin{itemize}
  \item[（1）] 证明：\(\beta, A \beta, \cdots, A^{n-1} \beta\) 线性无关。
  \item[（2）] 若 \(A^{n} \beta=A \beta\) ，求 \(A-I_{n}\) 的秩。
\end{itemize}

\subsection{题目6}
6．（18 分）设 \(n\) 阶矩阵 \(A\) 的特征多项式为 \(f(x)\) ，其中 \(f(x)=g(x) h(x),(g(x), h(x))=1\) ，且 \(\operatorname{deg}(g(x))<n, \operatorname{deg}(h(x))<n\) 。证明：存在可逆矩阵 \(P\) 和 \(m_{1}\) 阶矩阵 \(A_{1}, m_{2}\) 阶矩阵 \(A_{2}\) ，使得

\[
P^{-1} A P=\left(\begin{array}{ll}
A_{1} & \\
& A_{2}
\end{array}\right) .
\]

且 \(m_{1}, m_{2}<n\) ．

\subsection{题目7}
7．（18 分）设 \(A, B, X \in M^{3 \times 3}(\mathbb{C})\) ，对于复系数多项式 \(f(t), g(t)\) ，有 \(f(X)=A, g(X)=B\) ，且 \(A^{2}=B^{2}=O\) 。证明：存在 \(\lambda \in \mathbb{C}\) ，满足 \(A=\lambda B\) 或 \(B=\lambda A\) 。

\subsection{题目8}
8．（16 分）设 \(\varphi: M^{n \times n}(\mathbb{F}) \rightarrow M^{m \times m}(\mathbb{F})\) 为线性映射，且满足 \(\varphi\left(I_{n}\right)=I_{m}\) ，对任意的 \(A, B \in\) \(M^{n \times n}(\mathbb{F})\) ，有 \(\varphi(A B)=\varphi(A) \varphi(B)\) 。证明：

\begin{itemize}
  \item[（1）] \(\varphi\) 为单射．
  \item[（2）] \(n \mid m\) ．
\end{itemize}
\end{document}
